% =====================================================================
%  TP3 — Administration & Automatisation Azure (ShopEasy)
%  Binôme : Louis SCARFONE & Maxence BOURRAGUE — Bloc 4 (Cloud Computing)
%  Moteur : LuaLaTeX (LuaHBTeX)
% =====================================================================
\documentclass[11pt,a4paper]{article}

% ---------- Langue & fontes ----------
\usepackage{fontspec}
\usepackage{polyglossia}
\setdefaultlanguage{french}
\setmainfont{TeX Gyre Pagella}
\setsansfont{TeX Gyre Heros}
\setmonofont{DejaVu Sans Mono}[Scale=0.80]
\usepackage{microtype}
\emergencystretch=3em  % absorbe les débordements de texte (code inline en fin de ligne)

% ---------- Respiration (look épuré) ----------
\usepackage{setspace}
\setstretch{1.08}
\setlength{\parindent}{0pt}
\setlength{\parskip}{0.7em plus 0.12em minus 0.05em}

% ---------- Géométrie ----------
\usepackage[a4paper,margin=2.2cm,headheight=16pt,headsep=14pt,footskip=20pt]{geometry}

% ---------- Couleurs (palette Azure) ----------
\usepackage{xcolor}
\definecolor{azure}{HTML}{0078D4}
\definecolor{azuredark}{HTML}{14375E}
\definecolor{azuremid}{HTML}{2E75B6}
\definecolor{azurelight}{HTML}{E8F1FB}
\definecolor{azurepale}{HTML}{F3F8FD}
\definecolor{okgreen}{HTML}{107C10}
\definecolor{okgreenbg}{HTML}{E8F5E9}
\definecolor{warnorange}{HTML}{B5520A}
\definecolor{warnbg}{HTML}{FBEEDF}
\definecolor{dangerred}{HTML}{A4262C}
\definecolor{codebg}{HTML}{F6F8FA}
\definecolor{codeframe}{HTML}{D8E1EC}
\definecolor{ink}{HTML}{1B2733}
\definecolor{rowalt}{HTML}{EBF3FC}
\color{ink}

% ---------- Divers ----------
\usepackage{graphicx}
\graphicspath{{../livrables/assets/}}
\usepackage{fontawesome5}
\usepackage{enumitem}
\setlist{leftmargin=1.4em,topsep=2pt,itemsep=2pt,parsep=0pt}
\usepackage{needspace}
\usepackage{fvextra}

% ---------- Emojis (Twemoji couleur) ----------
\usepackage{twemojis}
\newcommand{\emo}[1]{\raisebox{-0.14em}{\twemoji{#1}}}
\newcommand{\eTarget}{\emo{1f3af}}
\newcommand{\ePackage}{\emo{1f4e6}}
\newcommand{\eCheck}{\emo{2705}}
\newcommand{\eWarn}{\emo{26a0}}
\newcommand{\eBulb}{\emo{1f4a1}}
\newcommand{\eLock}{\emo{1f512}}
\newcommand{\eKey}{\emo{1f510}}
\newcommand{\eRocket}{\emo{1f680}}
\newcommand{\eChart}{\emo{1f4ca}}
\newcommand{\eMoney}{\emo{1f4b0}}
\newcommand{\eGear}{\emo{2699}}
\newcommand{\eGlobe}{\emo{1f310}}
\newcommand{\eCamera}{\emo{1f4f8}}
\newcommand{\ePin}{\emo{1f4cc}}
\newcommand{\eMag}{\emo{1f50d}}
\newcommand{\eFlag}{\emo{1f3c1}}
\newcommand{\eParty}{\emo{1f389}}
\newcommand{\eShield}{\emo{1f6e1}}
\newcommand{\eBroom}{\emo{1f9f9}}
\newcommand{\eFolder}{\emo{1f4c2}}
\newcommand{\eServer}{\emo{1f5a5}}
\newcommand{\eDB}{\emo{1f5c4}}
\newcommand{\eScale}{\emo{2696}}
\newcommand{\eRed}{\emo{1f534}}
\newcommand{\eOrange}{\emo{1f7e0}}
\newcommand{\eYellow}{\emo{1f7e1}}
\newcommand{\eGreen}{\emo{1f7e2}}
\newcommand{\eWorld}{\emo{1f30d}}
\newcommand{\eTerminal}{\emo{1f4bb}}
\newcommand{\eSnake}{\emo{1f40d}}
\newcommand{\arr}{\ensuremath{\rightarrow}}
\newcommand{\code}[1]{\texttt{#1}}

% ---------- Tableaux (tabularray) ----------
\usepackage{tabularray}
\UseTblrLibrary{booktabs}
\NewDocumentEnvironment{ltable}{m}{%
  \par\addvspace{4pt}\begin{tblr}{
    width=\linewidth,
    colspec={#1},
    row{1}={bg=azuredark,fg=white,font=\sffamily\bfseries\small},
    row{2-Z}={font=\small},
    row{even}={bg=rowalt},
    hline{1,Z}={1.1pt,solid,azuredark},
    hline{2}={0.5pt,solid,azuredark!45},
    colsep=7pt, rowsep=4pt,
    stretch=1.18,
  }%
}{%
  \end{tblr}\par\addvspace{4pt}%
}

% ---------- Boîtes (tcolorbox) ----------
\usepackage[most]{tcolorbox}
\tcbset{
  boxbase/.style={
    enhanced, breakable, sharp corners, boxrule=0pt, leftrule=4pt,
    left=4.5mm, right=4.5mm, top=3.6mm, bottom=3.6mm,
    before skip=12pt, after skip=12pt,
    fonttitle=\sffamily\bfseries, coltitle=white,
    attach boxed title to top left={xshift=0mm,yshift=-2.6mm},
    boxed title style={sharp corners, boxrule=0pt, left=2.5mm,right=2.5mm,top=0.9mm,bottom=0.9mm},
  },
}
\newtcolorbox{fiche}{boxbase, colback=azurelight, colframe=azure}
\newtcolorbox{notebox}[1][Point de vigilance]{boxbase, colback=warnbg, colframe=warnorange,
  colbacktitle=warnorange, title={\eWarn~#1}}
\newtcolorbox{tipbox}[1][Astuce]{boxbase, colback=azurepale, colframe=azuremid,
  colbacktitle=azuremid, title={\eBulb~#1}}
\newtcolorbox{etatbox}[1][État après l'atelier]{boxbase, colback=okgreenbg, colframe=okgreen,
  colbacktitle=okgreen, title={\eCheck~#1}}

% Bloc code / terminal
\fvset{breaklines=true, breakanywhere=true, fontsize=\footnotesize,
  breaksymbolleft=\raisebox{0.1ex}{\scriptsize\textcolor{azure}{$\hookrightarrow$}}}
\newtcolorbox{codebox}[1]{
  enhanced, breakable, sharp corners, boxrule=0.7pt,
  colback=codebg, colframe=codeframe,
  coltitle=azuredark, colbacktitle=codeframe, fonttitle=\ttfamily\bfseries\scriptsize,
  title={\faTerminal~~#1}, left=3mm, right=3mm, top=2.6mm, bottom=2.6mm,
  before skip=12pt, after skip=14pt,
}

% Capture d'écran encadrée + légende
\newcommand{\screenshot}[2]{%
  \par\addvspace{10pt}
  \begin{center}
  \begin{tcolorbox}[enhanced, sharp corners, boxrule=0.7pt, colframe=codeframe, colback=white,
      left=0.8mm, right=0.8mm, top=0.8mm, bottom=0.8mm, drop fuzzy shadow=black!18]
    \includegraphics[width=\linewidth]{#1}
  \end{tcolorbox}\par\vspace{4pt}
  {\small\itshape\color{black!58}#2}
  \end{center}\par\addvspace{12pt}
}

% Carte « chiffre-clé » pour la page de garde
\newcommand{\statcard}[2]{%
  \begin{tcolorbox}[enhanced, sharp corners, boxrule=0pt, leftrule=3pt, colframe=azure,
     colback=azurepale, width=0.305\linewidth, halign=center, valign=center, height=2.3cm,
     left=1mm, right=1mm, top=1mm, bottom=1mm, nobeforeafter]
     {\sffamily\bfseries\fontsize{22}{24}\selectfont\color{azuredark}#1}\\[1.5mm]
     {\sffamily\small\color{black!58}#2}
  \end{tcolorbox}%
}

% ---------- Sections (titlesec) ----------
\usepackage{titlesec}
\titleformat{\section}
  {\sffamily\LARGE\bfseries\color{azuredark}}
  {\tcbox[on line, sharp corners, boxrule=0pt, colback=azure, colupper=white,
     left=2.2mm,right=2.2mm,top=0.6mm,bottom=0.6mm]{\thesection}}
  {0.7em}{}[\vspace{1mm}{\color{azure}\titlerule[1.6pt]}]
\titlespacing*{\section}{0pt}{3.6ex plus 0.8ex}{2.0ex}
\titleformat{\subsection}
  {\sffamily\large\bfseries\color{azuremid}}{\thesubsection}{0.6em}{}
\titlespacing*{\subsection}{0pt}{2.8ex plus 0.5ex}{1.1ex}
\titleformat{\subsubsection}
  {\sffamily\bfseries\color{azuredark}}{}{0pt}{}
\titlespacing*{\subsubsection}{0pt}{1.9ex plus 0.3ex}{0.5ex}

% ---------- En-têtes / pieds (fancyhdr) ----------
\usepackage{fancyhdr}
\pagestyle{fancy}
\fancyhf{}
\fancyhead[L]{\sffamily\small\color{azuredark}\textbf{TP3 Administration}~\textcolor{azure}{\textbullet}~ShopEasy}
\fancyhead[R]{\sffamily\small\color{black!55}\nouppercase{\leftmark}}
\fancyfoot[C]{\sffamily\small\color{black!55}\thepage}
\fancyfoot[R]{\sffamily\scriptsize\color{black!42}Bloc 4 · Cloud Computing}
\fancyfoot[L]{\sffamily\scriptsize\color{black!42}SCARFONE · BOURRAGUE}
\renewcommand{\headrulewidth}{0pt}
\renewcommand{\footrulewidth}{0pt}
\renewcommand{\headrule}{\vspace{-2pt}{\color{azure!70}\rule{\headwidth}{1.1pt}}}

% ---------- Liens ----------
\usepackage[hidelinks]{hyperref}
\hypersetup{colorlinks=true, linkcolor=azuredark, urlcolor=azuremid,
  pdftitle={TP3 Administration - ShopEasy}, pdfauthor={Louis SCARFONE, Maxence BOURRAGUE}}

% Filet décoratif
\newcommand{\thinrule}{\par\vspace{1mm}{\color{azure!30}\hrule height 0.5pt}\vspace{1mm}}

\begin{document}

% =====================================================================
%  PAGE DE GARDE
% =====================================================================
\begin{titlepage}
\thispagestyle{empty}
\setlength{\parskip}{0pt}
\raggedright

\begin{tcolorbox}[enhanced, sharp corners, boxrule=0pt, colback=azuredark,
   width=\linewidth, left=13mm, right=13mm, top=14mm, bottom=14mm,
   borderline north={5pt}{0pt}{azure}, borderline south={2pt}{0pt}{azure}]
  \raggedright
  {\sffamily\bfseries\large\color{azure!75}\faMicrosoft~~MICROSOFT AZURE \quad\textbullet\quad BLOC 4 — CLOUD COMPUTING}

  \vspace{9mm}
  {\sffamily\bfseries\color{white}\fontsize{36}{42}\selectfont TP3 — Administration \& Automatisation}

  \vspace{4mm}
  {\sffamily\color{azurelight}\LARGE exploiter l'environnement Azure de \textbf{ShopEasy}}

  \vspace{8mm}
  {\color{azure}\rule{5.2cm}{2.5pt}}

  \vspace{5mm}
  {\sffamily\color{white}\large Inventorier, administrer, superviser et optimiser une infrastructure Azure avec \textbf{Azure CLI}, \textbf{Bash}, \textbf{Python} et \textbf{Azure Monitor}.}
\end{tcolorbox}

\vspace{1.8cm}

\begin{minipage}[t]{0.5\linewidth}
  {\sffamily\bfseries\large\color{azuredark}Binôme}\\[3.5mm]
  {\large Louis \textsc{Scarfone}}\\[2mm]
  {\large Maxence \textsc{Bourragué}}
\end{minipage}\hfill
\begin{minipage}[t]{0.46\linewidth}
  {\sffamily\bfseries\large\color{azuredark}Cadre}\\[3.5mm]
  Mastère Dev Manager Full Stack — RNCP niveau 7\\[2mm]
  Cas fil rouge : \textbf{ShopEasy}\\[2mm]
  Azure for Students — région \texttt{swedencentral}
\end{minipage}

\vspace{1.9cm}

\noindent
\statcard{12}{ateliers + quiz}\hfill
\statcard{4}{scripts CLI/Bash/Python}\hfill
\statcard{$\approx$\,47\,\$}{coût mensuel estimé}

\vfill
{\sffamily\large\bfseries\color{azuredark}Juin 2026}\hfill
{\sffamily\color{black!45}Mastère Dev Manager Full Stack}
\end{titlepage}

% =====================================================================
%  TABLE DES MATIÈRES
% =====================================================================
\thispagestyle{fancy}
{\sffamily\Huge\bfseries\color{azuredark}Sommaire}
\vspace{2mm}{\color{azure}\titlerule[1.6pt]}\vspace{4mm}
\hypersetup{linkcolor=azuredark}
{\setlength{\parskip}{1pt}\tableofcontents}
\newpage


% #####################################################################
%  ATELIER 1
% #####################################################################
\section{Atelier 1 — Prise en main d'Azure CLI}

\begin{fiche}
\textbf{\eTarget~Objectif} — vérifier l'accès à l'environnement Azure, identifier la subscription active et se positionner sur le bon groupe de ressources.

\textbf{\ePackage~Livrable attendu} — extrait de terminal montrant la subscription active, le groupe de ressources cible et les valeurs par défaut configurées.
\end{fiche}

\subsection{Contexte — du déploiement à l'exploitation}
Le TP1 a \textbf{conçu et déployé manuellement} l'architecture ShopEasy ; le TP2 l'a \textbf{décrite en Terraform}. Le TP3 ouvre la phase d'\textbf{exploitation} : administrer, inventorier, surveiller, automatiser et documenter cet environnement avec \textbf{Azure CLI, Bash et Python}. L'environnement de travail est celui du TP2, \textbf{redéployé à l'identique} via \code{terraform apply} (21 ressources) : un groupe de ressources \code{rg-shopeasy-dev} contenant un réseau segmenté, deux VM Linux derrière un Load Balancer et un compte de stockage privé.

\begin{ltable}{Q[l,0.5] X[l,1]}
Périmètre & Valeur \\
Abonnement & Azure for Students \\
Région & \code{swedencentral} (Stockholm) \\
Groupe de ressources & \code{rg-shopeasy-dev} \\
VM web & \code{vm-shopeasy-dev-web-1}, \code{vm-shopeasy-dev-web-2} (\code{Standard\_B2ats\_v2}) \\
Storage Account & \code{shopeasydevdocs350wnq} \\
\end{ltable}

\begin{tipbox}[Contraintes Azure for Students]
La région \code{swedencentral} et la taille \code{Standard\_B2ats\_v2} sont conservées du TP1/TP2 : \code{francecentral} est bloquée par la policy \textit{Allowed regions}, et \code{Standard\_B1s} est indisponible sur l'abonnement étudiant.
\end{tipbox}

\subsection{Centralisation des variables — \code{variables.sh}}
Toutes les commandes et scripts s'appuient sur un jeu de variables centralisé dans \code{variables.sh}, à sourcer avant chaque session. Cette factorisation évite de répéter les noms de ressources et permet de rejouer les commandes sur un autre environnement.

\begin{codebox}{variables.sh}
\begin{Verbatim}
#!/usr/bin/env bash
# variables.sh - Variables d'exploitation de l'environnement ShopEasy (TP3)

export RG="rg-shopeasy-dev"              # Groupe de ressources dedie au projet
export LOCATION="swedencentral"          # Region de deploiement
export VM1="vm-shopeasy-dev-web-1"       # Premiere VM web
export VM2="vm-shopeasy-dev-web-2"       # Seconde VM web
export STORAGE="shopeasydevdocs350wnq"   # Storage Account (suffixe aleatoire Terraform)
export CONTAINER="operations"            # Conteneur des rapports d'exploitation

export ARM_SUBSCRIPTION_ID="$(az account show --query id -o tsv)"
export AZURE_SUBSCRIPTION_ID="$ARM_SUBSCRIPTION_ID"   # Utilisee par le SDK Python
export AZURE_RESOURCE_GROUP="$RG"
\end{Verbatim}
\end{codebox}

Aucun identifiant n'est codé en dur : la subscription est résolue dynamiquement par \code{az account show}.

\subsection{Version d'Azure CLI \& subscription active}
\begin{codebox}{bash}
\begin{Verbatim}
az version
az account show --output table
az account list --output table
\end{Verbatim}
\end{codebox}
\begin{codebox}{console}
\begin{Verbatim}
azure-cli : 2.87.0

EnvironmentName  IsDefault  Name                State    TenantDefaultDomain  TenantId
---------------  ---------  ------------------  -------  -------------------  --------
AzureCloud       True       Azure for Students  Enabled  efrei.net            <masque>
\end{Verbatim}
\end{codebox}
La subscription \textbf{Azure for Students} est active (\code{State = Enabled}) et définie par défaut (\code{IsDefault = True}). Les identifiants (subscription ID, tenant ID) sont \textbf{masqués} dans les livrables ; ils ne sont jamais écrits en clair dans un fichier versionné.

\subsection{Groupe de ressources cible}
\begin{codebox}{bash}
\begin{Verbatim}
az group show --name "$RG" \
  --query "{Nom:name, Region:location, Etat:properties.provisioningState}" -o table
\end{Verbatim}
\end{codebox}
\begin{codebox}{console}
\begin{Verbatim}
Nom              Region         Etat
---------------  -------------  ---------
rg-shopeasy-dev  swedencentral  Succeeded
\end{Verbatim}
\end{codebox}
Le groupe \code{rg-shopeasy-dev} existe, en région \code{swedencentral}, à l'état \code{Succeeded} : l'environnement à exploiter est bien en place.

\subsection{Configuration des valeurs par défaut}
Pour éviter de répéter \code{--resource-group} et \code{--location} sur chaque commande, on configure des valeurs par défaut au niveau de la CLI.
\begin{codebox}{bash}
\begin{Verbatim}
az configure --defaults group="$RG" location="$LOCATION"
az configure --list-defaults --output table
\end{Verbatim}
\end{codebox}
\begin{codebox}{console}
\begin{Verbatim}
Name      Source            Value
--------  ----------------  ---------------
group     ~/.azure/config   rg-shopeasy-dev
location  ~/.azure/config   swedencentral
\end{Verbatim}
\end{codebox}
Effet immédiat : les commandes suivantes ciblent le bon groupe \textbf{sans \code{-g}} ; \code{az resource list} renvoie alors directement les \textbf{13 ressources} de premier niveau.

\subsection{Travail demandé}
\begin{enumerate}
  \item \textbf{Vérifier la subscription active.} \code{az account show} \arr{} \textbf{Azure for Students}, \textit{Enabled}, par défaut.
  \item \textbf{Vérifier que le groupe de ressources existe.} \code{az group show} \arr{} \code{rg-shopeasy-dev}, \code{swedencentral}, \code{Succeeded}.
  \item \textbf{Configurer les valeurs par défaut.} \code{az configure --defaults} \arr{} groupe et région enregistrés dans \code{\textasciitilde/.azure/config}, commandes ultérieures raccourcies.
\end{enumerate}

\subsection{Capture}
\screenshot{atelier01-resource-group.png}{Groupe de ressources \code{rg-shopeasy-dev} en région \textit{Sweden Central} avec ses ressources (Portal \arr{} Resource groups \arr{} rg-shopeasy-dev).}

\begin{etatbox}
\begin{itemize}
  \item Accès Azure validé : subscription \textbf{Azure for Students} active (\textit{Enabled}, par défaut).
  \item Groupe \code{rg-shopeasy-dev} présent en \code{swedencentral} (\code{Succeeded}), contenant les 13 ressources de l'architecture.
  \item Valeurs par défaut configurées (\code{group}, \code{location}) : commandes plus courtes et reproductibles.
  \item Variables d'exploitation centralisées dans \code{variables.sh} (aucun secret en dur).
  \item \textbf{Prêt pour l'Atelier 2 — Inventorier les ressources Azure.}
\end{itemize}
\end{etatbox}


% #####################################################################
%  ATELIER 2
% #####################################################################
\section{Atelier 2 — Inventorier les ressources Azure}

\begin{fiche}
\textbf{\eTarget~Objectif} — produire un inventaire lisible des ressources ShopEasy, rejouable sans passer par le portail.

\textbf{\ePackage~Livrable attendu} — fichiers \code{resources.json}, \code{resources.tsv} et tableau d'inventaire complété.
\end{fiche}

\subsection{Lister les ressources}
L'inventaire est la base de l'exploitation : avant de surveiller, sécuriser ou optimiser, il faut savoir ce qui existe. Azure CLI liste les ressources et choisit le format de sortie selon l'usage.
\begin{codebox}{bash}
\begin{Verbatim}
az resource list --resource-group "$RG" \
  --query "[].{Nom:name,Type:type,Region:location}" --output table
\end{Verbatim}
\end{codebox}
\begin{codebox}{console}
\begin{Verbatim}
Nom                     Type                                     Region
----------------------  ---------------------------------------  -------------
vnet-shopeasy-dev       Microsoft.Network/virtualNetworks        swedencentral
shopeasydevdocs350wnq   Microsoft.Storage/storageAccounts        swedencentral
pip-shopeasy-dev-web-1  Microsoft.Network/publicIPAddresses      swedencentral
pip-shopeasy-dev-lb     Microsoft.Network/publicIPAddresses      swedencentral
nsg-shopeasy-dev-web    Microsoft.Network/networkSecurityGroups  swedencentral
pip-shopeasy-dev-web-2  Microsoft.Network/publicIPAddresses      swedencentral
lb-shopeasy-dev-web     Microsoft.Network/loadBalancers          swedencentral
nic-shopeasy-dev-web-2  Microsoft.Network/networkInterfaces      swedencentral
nic-shopeasy-dev-web-1  Microsoft.Network/networkInterfaces      swedencentral
vm-shopeasy-dev-web-2   Microsoft.Compute/virtualMachines        swedencentral
vm-shopeasy-dev-web-1   Microsoft.Compute/virtualMachines        swedencentral
..._OsDisk_... (x2)     Microsoft.Compute/disks                  swedencentral
\end{Verbatim}
\end{codebox}

\subsection{Exporter l'inventaire (JSON + TSV)}
Un inventaire d'exploitation doit être \textbf{conservé et rejouable}. On exporte le \textbf{JSON} complet (archivage, traitement) et le \textbf{TSV} (réutilisation tableur/script).
\begin{codebox}{bash}
\begin{Verbatim}
az resource list --resource-group "$RG" --output json > exports/resources.json
az resource list --resource-group "$RG" \
  --query "[].{Nom:name,Type:type,Region:location}" --output tsv > exports/resources.tsv
\end{Verbatim}
\end{codebox}
Les ressources portent déjà des \textbf{tags hérités du TP2} (appliqués par Terraform, en \code{snake\_case}). L'Atelier 3 normalisera une stratégie de tags d'exploitation.

\subsection{Identifier les types}
\begin{codebox}{bash}
\begin{Verbatim}
az resource list --resource-group "$RG" --query "[].type" -o tsv | sort | uniq -c | sort -rn
\end{Verbatim}
\end{codebox}
\begin{codebox}{console}
\begin{Verbatim}
   3 Microsoft.Network/publicIPAddresses
   2 Microsoft.Network/networkInterfaces
   2 Microsoft.Compute/virtualMachines
   2 Microsoft.Compute/disks
   1 Microsoft.Storage/storageAccounts
   1 Microsoft.Network/virtualNetworks
   1 Microsoft.Network/networkSecurityGroups
   1 Microsoft.Network/loadBalancers
\end{Verbatim}
\end{codebox}
\textbf{8 types} pour \textbf{13 ressources} de premier niveau.

\begin{notebox}[Ressources de premier niveau seulement]
\code{az resource list} ne renvoie que les ressources de premier niveau. Les sous-ressources (subnet, règles NSG, règle/sonde/backend pool du Load Balancer) sont imbriquées dans leur parent. Côté Terraform, l'état gère \textbf{21 objets}.
\end{notebox}

\subsection{Tableau d'inventaire}
\begin{ltable}{Q[l,0.62] Q[l,0.5] X[l,1]}
Ressource & Type & Rôle dans l'architecture \\
\code{vnet-shopeasy-dev} & VNet & Réseau privé \code{10.20.0.0/16}, segmentation (subnet \code{snet-web}). \\
\code{nsg-shopeasy-dev-web} & NSG & Pare-feu du subnet : HTTP (80) ouvert, SSH (22) restreint à l'IP admin. \\
\code{lb-shopeasy-dev-web} & Load Balancer & LB \textbf{Standard} L4 : répartit le trafic HTTP sur les 2 VM. \\
\code{pip-shopeasy-dev-lb} & Public IP & IP publique du LB = point d'entrée HTTP unique. \\
\code{pip-...-web-1/2} & Public IP & IP publiques d'administration des 2 VM. \\
\code{nic-...-web-1/2} & NIC & Cartes réseau des 2 VM (subnet + backend pool). \\
\code{vm-...-web-1/2} & VM & Serveurs web : Ubuntu 22.04 + Nginx (\code{Standard\_B2ats\_v2}). \\
\code{...\_OsDisk\_...} ($\times$2) & Disk & Disques OS managés des 2 VM. \\
\code{shopeasydevdocs350wnq} & Storage & Stockage documentaire privé (\code{StorageV2}, LRS, versioning). \\
\end{ltable}

\subsection{Travail demandé}
\begin{enumerate}
  \item \textbf{Tableau d'inventaire (nom, type, région, rôle).} Réalisé : 13 ressources, toutes en \code{swedencentral}, reliées à leur rôle.
  \item \textbf{Inventaire rejouable sans le portail.} Produit en CLI, exporté en \code{resources.json} et \code{resources.tsv}, industrialisé ensuite dans \code{inventory.sh} (Atelier 5).
\end{enumerate}

\begin{etatbox}
\begin{itemize}
  \item Inventaire CLI complet : \textbf{13 ressources}, \textbf{8 types}, région unique \code{swedencentral}.
  \item Exports générés : \code{resources.json} (368 lignes) et \code{resources.tsv} (13 lignes).
  \item Tableau d'inventaire rôle par rôle établi (base du rapport d'exploitation).
  \item Constat : ressources déjà taguées (TP2), à normaliser à l'Atelier 3.
  \item \textbf{Prêt pour l'Atelier 3 — Normaliser les tags d'exploitation.}
\end{itemize}
\end{etatbox}


% #####################################################################
%  ATELIER 3
% #####################################################################
\section{Atelier 3 — Normaliser les tags d'exploitation}

\begin{fiche}
\textbf{\eTarget~Objectif} — appliquer une stratégie de tags minimale pour faciliter l'exploitation, le suivi des coûts et l'identification des responsabilités.

\textbf{\ePackage~Livrable attendu} — export JSON (et capture) montrant les tags du groupe de ressources et d'au moins une VM.
\end{fiche}

\subsection{Stratégie de tags d'exploitation}
Les ressources avaient déjà des \textbf{tags techniques hérités du TP2} (posés par Terraform, en \code{snake\_case}). On les \textbf{normalise} en une convention d'exploitation homogène (\code{PascalCase}).
\begin{ltable}{Q[l,0.4] Q[l,0.6] X[l,1]}
Tag & Valeur & Rôle \\
\code{Application} & \code{shopeasy} & Relier la ressource à l'application métier. \\
\code{Environment} & \code{dev} & Distinguer dev / test / prod. \\
\code{Owner} & \code{cloudops-shopeasy} & Identifier l'équipe responsable. \\
\code{CostCenter} & \code{cloud-training} & Imputer les coûts (FinOps). \\
\code{ManagedBy} & \code{mixed} & Provisionné par Terraform, exploité par CLI. \\
\code{Role} (VM) & \code{web} & Rôle technique de la VM. \\
\end{ltable}

\begin{tipbox}[Choix ManagedBy = mixed]
L'infrastructure a été créée en IaC (Terraform, TP2) mais est administrée en CLI au TP3. La valeur \code{mixed} décrit honnêtement ce cycle de vie hybride, plutôt que \code{cli} (qui masquerait l'origine IaC) ou \code{terraform} (qui ignorerait l'exploitation CLI).
\end{tipbox}

\subsection{État avant normalisation}
Le RG et les VM portaient les 5 tags Terraform en \code{snake\_case} :
\begin{codebox}{console}
\begin{Verbatim}
{
  "cost_center": "cloud-training",
  "environment": "dev",
  "managed_by": "terraform",
  "owner": "formation",
  "project": "shopeasy"
}
\end{Verbatim}
\end{codebox}

\subsection{Appliquer les tags (RG puis VM)}
\code{az group update --tags} \textbf{remplace} l'intégralité du jeu de tags du RG. Pour les VM, \code{az resource tag --ids} remplace proprement et ajoute \code{Role=web}.
\begin{codebox}{bash}
\begin{Verbatim}
az group update --name "$RG" \
  --tags Application=shopeasy Environment=dev Owner=cloudops-shopeasy \
         CostCenter=cloud-training ManagedBy=mixed

for VM in "$VM1" "$VM2"; do
  VM_ID=$(az vm show -g "$RG" -n "$VM" --query id -o tsv)
  az resource tag --ids "$VM_ID" \
    --tags Application=shopeasy Environment=dev Owner=cloudops-shopeasy \
           CostCenter=cloud-training ManagedBy=mixed Role=web
done
\end{Verbatim}
\end{codebox}
\begin{codebox}{console}
\begin{Verbatim}
Application    Environment    Owner              CostCenter      ManagedBy
-------------  -------------  -----------------  --------------  -----------
shopeasy       dev            cloudops-shopeasy  cloud-training  mixed
\end{Verbatim}
\end{codebox}

\subsection{Vérification et couverture}
Le travail demandé porte sur le RG et les VM. Les autres ressources conservent les tags techniques Terraform :
\begin{codebox}{bash}
\begin{Verbatim}
az resource list -g "$RG" --query "[?tags.Application==null].name" -o tsv
\end{Verbatim}
\end{codebox}
\textbf{11 ressources} ne portent pas encore le tag \code{Application} normalisé. Ce manque est \textbf{volontairement laissé visible} : il est \textbf{détecté automatiquement} par \code{inventory.sh} (Atelier 5) et \code{healthcheck.sh} (Atelier 9), et son extension figure dans les recommandations FinOps (Atelier 11).

\begin{notebox}[Drift Terraform]
Modifier les tags via CLI sur des ressources gérées par Terraform crée une \textbf{dérive (drift)} : au prochain \code{terraform plan}, Terraform proposerait de restaurer ses tags \code{snake\_case} (phénomène étudié au TP2, Atelier 11). Réconciliation : reporter la convention dans le code (\code{locals.common\_tags}) ou exclure les tags (\code{lifecycle ignore\_changes}). Pour le TP3, aucun \code{apply} n'est rejoué d'ici la destruction finale.
\end{notebox}

\subsection{Pourquoi les tags sont utiles (FinOps \& gouvernance)}
\textbf{FinOps} — les tags \code{CostCenter}, \code{Application} et \code{Environment} permettent de \textbf{ventiler la facture} Azure (Cost Management) par centre de coût, application et environnement. Sans eux, impossible de répondre à « combien coûte ShopEasy en dev ? ». \textbf{Gouvernance} — \code{Owner} identifie le responsable (incident, autorisation d'arrêt), \code{ManagedBy} documente le mode de gestion ; les tags servent de base à des \textbf{Azure Policy} et au filtrage des inventaires.

\subsection{Captures}
\screenshot{atelier03-tags-rg.png}{Tags normalisés du groupe \code{rg-shopeasy-dev} (Portal \arr{} Resource groups \arr{} rg-shopeasy-dev \arr{} Tags) : les 5 tags d'exploitation.}
\screenshot{atelier03-tags-vm.png}{Tags normalisés de la VM \code{vm-shopeasy-dev-web-1} (Portal \arr{} Virtual machines \arr{} \dots{} \arr{} Tags) : les 6 tags dont \code{Role=web}.}

\begin{etatbox}
\begin{itemize}
  \item Convention normalisée (\code{Application}, \code{Environment}, \code{Owner}, \code{CostCenter}, \code{ManagedBy}) appliquée au \textbf{RG} et aux \textbf{2 VM} (+ \code{Role=web}).
  \item Anciennes clés \code{snake\_case} remplacées proprement ; preuve exportée (\code{tags-rg.json}, \code{tags-vm.json}).
  \item \textbf{11 ressources} encore en \code{snake\_case} : manque tracé, à corriger via inventaire automatisé.
  \item Drift Terraform identifié et documenté (réconciliation par le code ou \code{ignore\_changes}).
  \item \textbf{Prêt pour l'Atelier 4 — Administrer les machines virtuelles.}
\end{itemize}
\end{etatbox}


% #####################################################################
%  ATELIER 4
% #####################################################################
\section{Atelier 4 — Administrer les machines virtuelles}

\begin{fiche}
\textbf{\eTarget~Objectif} — réaliser les opérations courantes sur les VM : lister, vérifier l'état, démarrer, arrêter, désallouer, redémarrer et exécuter une commande distante.

\textbf{\ePackage~Livrable attendu} — tableau des VM (état, IP, taille, action), avec la commande utilisée et le résultat observé.
\end{fiche}

\subsection{Lister les VM et leur état}
\code{--show-details} ajoute l'état d'alimentation et l'IP publique.
\begin{codebox}{bash}
\begin{Verbatim}
az vm list -g "$RG" --show-details \
  --query "[].{Nom:name,Etat:powerState,IP:publicIps,Taille:hardwareProfile.vmSize}" -o table
\end{Verbatim}
\end{codebox}
\begin{codebox}{console}
\begin{Verbatim}
Nom                    Etat        IP              Taille
---------------------  ----------  --------------  -----------------
vm-shopeasy-dev-web-1  VM running  20.240.233.166  Standard_B2ats_v2
vm-shopeasy-dev-web-2  VM running  20.91.246.70    Standard_B2ats_v2
\end{Verbatim}
\end{codebox}

\subsection{Exécuter une commande distante (vérifier le serveur web)}
\code{az vm run-command invoke} exécute un script dans la VM via l'agent Azure, \textbf{sans connexion SSH} (utile quand le port 22 est restreint à l'IP admin).
\begin{codebox}{bash}
\begin{Verbatim}
az vm run-command invoke -g "$RG" -n "$VM1" --command-id RunShellScript \
  --scripts "hostname && uptime && systemctl is-active nginx && \
             curl -s -o /dev/null -w 'HTTP local: %{http_code}\n' http://localhost" \
  --query "value[0].message" -o tsv
\end{Verbatim}
\end{codebox}
\begin{codebox}{console}
\begin{Verbatim}
[stdout]
vm-shopeasy-dev-web-1
 08:15:51 up 21 min,  0 users,  load average: 0.00, 0.00, 0.00
active
nginx.service - active (running) since Fri 2026-06-26 07:55:02 UTC
HTTP local: 200
\end{Verbatim}
\end{codebox}
\textbf{Lecture} : le serveur répond, \code{nginx} est \code{active (running)} et la requête HTTP locale renvoie \textbf{200}. Le service web est opérationnel.

\subsection{Redémarrer et désallouer}
Deux actions d'exploitation sur deux VM distinctes : \textbf{redémarrer} la VM 1 (reste allouée) et \textbf{désallouer} la VM 2 non utilisée (geste FinOps : libère le calcul).
\begin{codebox}{bash}
\begin{Verbatim}
az vm restart    -g "$RG" -n "$VM1"   # redemarrage, VM toujours allouee
az vm deallocate -g "$RG" -n "$VM2"   # arret + liberation du compute
\end{Verbatim}
\end{codebox}
\begin{codebox}{console}
\begin{Verbatim}
Nom                    Etat            IP              Taille
---------------------  --------------  --------------  -----------------
vm-shopeasy-dev-web-1  VM running      20.240.233.166  Standard_B2ats_v2
vm-shopeasy-dev-web-2  VM deallocated  20.91.246.70    Standard_B2ats_v2
\end{Verbatim}
\end{codebox}

\subsection{\code{stop} $\neq$ \code{deallocate} (point FinOps essentiel)}
\begin{ltable}{Q[l,0.5] X[l,1] Q[c,0.5]}
Commande & Effet & Compute facturé \\
\code{az vm stop} & Arrête l'OS, mais la VM \textbf{reste allouée} & \textbf{Continue} \\
\code{az vm deallocate} & Arrête \textbf{et libère} le calcul (\textit{Stopped (deallocated)}) & \textbf{Stoppé} \\
\end{ltable}
\begin{notebox}[Ce qui reste facturé même désallouée]
Le \textbf{disque OS} et l'\textbf{IP publique Standard statique} (réservée) restent facturés : \code{vm-shopeasy-dev-web-2} affiche encore son IP bien qu'éteinte. La suppression des IP publiques (Bastion) est un levier FinOps de l'Atelier 11.
\end{notebox}

\subsection{Tableau récapitulatif des VM}
\begin{ltable}{Q[l,0.62] Q[l,0.45] Q[l,0.5] X[l,0.9]}
VM & Taille & État final & Action \& résultat \\
\code{vm-...-web-1} & \code{B2ats\_v2} & \textbf{running} & \code{run-command} (Nginx \code{active}, HTTP 200) + \code{restart} \\
\code{vm-...-web-2} & \code{B2ats\_v2} & \textbf{deallocated} & \code{deallocate} (compute libéré, FinOps) \\
\end{ltable}

\subsection{Travail demandé}
\begin{enumerate}
  \item \textbf{Lister toutes les VM.} \code{az vm list} \arr{} 2 VM.
  \item \textbf{Identifier taille, IP, état.} \code{--show-details} \arr{} \code{B2ats\_v2}, IP publiques, \code{VM running}.
  \item \textbf{Redémarrer une VM.} \code{az vm restart} \arr{} reste \code{running}.
  \item \textbf{Désallouer une VM non utilisée.} \code{az vm deallocate} \arr{} \code{VM deallocated}.
  \item \textbf{Commande distante.} \code{az vm run-command invoke} \arr{} Nginx \code{active}, HTTP 200.
\end{enumerate}

\subsection{Capture}
\screenshot{atelier04-vms-states.png}{Liste des VM (Portal \arr{} Virtual machines) : \code{web-1} \textit{Running} et \code{web-2} \textit{Stopped (deallocated)} après les actions.}

\begin{etatbox}
\begin{itemize}
  \item 2 VM \code{Standard\_B2ats\_v2} inventoriées (état, IP, taille).
  \item Commande distante validée (sans SSH) : \code{web-1} sert Nginx, HTTP 200.
  \item \code{restart} sur la VM 1 (\code{running}), \code{deallocate} sur la VM 2 (\code{deallocated}, compute libéré).
  \item Distinction \code{stop} / \code{deallocate} documentée (impact FinOps).
  \item \textbf{Prêt pour l'Atelier 5 — Construire un script Bash d'inventaire.}
\end{itemize}
\end{etatbox}


% #####################################################################
%  ATELIER 5
% #####################################################################
\section{Atelier 5 — Construire un script Bash d'inventaire}

\begin{fiche}
\textbf{\eTarget~Objectif} — transformer les commandes manuelles d'inventaire en script réutilisable.

\textbf{\ePackage~Livrable attendu} — script \code{inventory.sh} fonctionnel, captures d'exécution et fichiers générés dans \code{exports/}.
\end{fiche}

\subsection{Le script \code{scripts/inventory.sh}}
Le script reprend les commandes des Ateliers 2 à 4 et y ajoute une \textbf{synthèse chiffrée} et des \textbf{contrôles d'exploitation}. Paramétrable, robuste (\code{set -euo pipefail}), exports \textbf{horodatés}.
\begin{codebox}{scripts/inventory.sh}
\begin{Verbatim}
#!/usr/bin/env bash
set -euo pipefail

RG="${1:-rg-shopeasy-dev}"          # groupe de ressources (argument ou defaut)
DATE="$(date +%Y%m%d-%H%M%S)"       # horodatage des fichiers d'export
OUT_DIR="exports"; mkdir -p "$OUT_DIR"

# 1. Verification de l'existence du groupe de ressources
if ! az group show --name "$RG" >/dev/null 2>&1; then
  echo "ERREUR : le groupe de ressources '$RG' est introuvable." >&2; exit 1
fi

# 2. Export JSON des ressources
az resource list --resource-group "$RG" \
  --query "[].{name:name,type:type,location:location,tags:tags,id:id}" \
  --output json > "$OUT_DIR/resources-$DATE.json"

# 3. Export TXT de l'etat des VM
az vm list --resource-group "$RG" --show-details \
  --query "[].{name:name,powerState:powerState,publicIps:publicIps,vmSize:hardwareProfile.vmSize}" \
  --output table > "$OUT_DIR/vms-$DATE.txt"

# 4. Synthese chiffree
TOTAL=$(az resource list -g "$RG" --query "length(@)" -o tsv)
NB_VM=$(az vm list -g "$RG" --query "length(@)" -o tsv)
NB_STORAGE=$(az storage account list -g "$RG" --query "length(@)" -o tsv)
NB_UNTAGGED=$(az resource list -g "$RG" --query "length([?tags.Application==null])" -o tsv)
printf "  %-34s %s\n" "Ressources (total)"              "$TOTAL"
printf "  %-34s %s\n" "Machines virtuelles"             "$NB_VM"
printf "  %-34s %s\n" "Comptes de stockage"             "$NB_STORAGE"
printf "  %-34s %s\n" "Ressources sans tag Application" "$NB_UNTAGGED"

# 5. Controles : VM en cours d'execution (compute facture)
RUNNING=$(az vm list -g "$RG" --show-details \
            --query "[?powerState=='VM running'].name" -o tsv)
if [[ -n "$RUNNING" ]]; then
  echo "  [AVERTISSEMENT] VM en cours d'execution (compute facture) :"
  echo "$RUNNING" | sed 's/^/    - /'
fi
if [[ "$NB_UNTAGGED" -gt 0 ]]; then
  echo "  [INFO] $NB_UNTAGGED ressource(s) sans tag 'Application'"
fi
\end{Verbatim}
\end{codebox}

\subsection{Exécution réelle}
\begin{codebox}{console}
\begin{Verbatim}
[1/4] Export des ressources -> exports/resources-20260626-103147.json
[2/4] Export de l'etat des VM -> exports/vms-20260626-103147.txt
[3/4] Calcul de la synthese...

--- Synthese ------------------------------------------------
  Ressources (total)                 13
  Machines virtuelles                2
  Comptes de stockage                1
  Ressources sans tag Application    11
-------------------------------------------------------------
[4/4] Controles d'exploitation...

  [AVERTISSEMENT] VM en cours d'execution (compute facture) :
    - vm-shopeasy-dev-web-1
  [INFO] 11 ressource(s) sans tag 'Application'
\end{Verbatim}
\end{codebox}

\begin{tipbox}[Validation croisée]
L'avertissement ne liste \textbf{que} \code{vm-shopeasy-dev-web-1} (en \code{VM running}), et \textbf{pas} \code{vm-shopeasy-dev-web-2} (désallouée à l'Atelier 4). Le filtre \code{powerState=='VM running'} fonctionne donc correctement.
\end{tipbox}

\subsection{Améliorations apportées (travail demandé)}
\begin{ltable}{Q[l,1] Q[l,1] Q[c,0.35]}
Amélioration demandée & Implémentation (JMESPath) & Résultat \\
Nombre total de ressources & \code{length(@)} & \textbf{13} \\
Nombre de VM & \code{az vm list ... length(@)} & \textbf{2} \\
Nombre de comptes de stockage & \code{az storage account list ... length(@)} & \textbf{1} \\
Ressources sans tag \code{Application} & \code{length([?tags.Application==null])} & \textbf{11} \\
Avertissement si VM \code{running} & \code{[?powerState=='VM running']} & web-1 signalée \\
\end{ltable}

\subsection{Bonnes pratiques Bash}
\begin{ltable}{Q[l,0.5] X[l,1]}
Bonne pratique & Mise en œuvre \\
Arrêt sur erreur & \code{set -euo pipefail} (erreur, variable non définie, échec dans un pipe). \\
Paramétrage & \code{RG="\$\{1:-rg-shopeasy-dev\}"} : groupe en argument, valeur par défaut. \\
Vérification préalable & \code{az group show ... || exit 1} : sortie propre si le RG n'existe pas. \\
Traçabilité & Exports \textbf{horodatés} (\code{-\$DATE}) : pas d'écrasement, historique conservé. \\
Idempotence & \code{mkdir -p}, relances sans effet de bord. \\
\end{ltable}

\begin{etatbox}
\begin{itemize}
  \item Script \code{scripts/inventory.sh} fonctionnel, commenté, robuste et paramétrable.
  \item Synthèse automatisée : \textbf{13 ressources, 2 VM, 1 compte de stockage, 11 sans tag \code{Application}}.
  \item Avertissement FinOps opérationnel (VM \code{running} détectée), filtrage d'état validé.
  \item Exports horodatés générés dans \code{exports/} (JSON ressources + TXT VM).
  \item \textbf{Prêt pour l'Atelier 6 — Automatiser l'arrêt et le démarrage des VM.}
\end{itemize}
\end{etatbox}


% #####################################################################
%  ATELIER 6
% #####################################################################
\section{Atelier 6 — Automatiser l'arrêt et le démarrage des VM}

\begin{fiche}
\textbf{\eTarget~Objectif} — construire un script pilotant l'état des VM (start / stop / deallocate / status) de manière sûre.

\textbf{\ePackage~Livrable attendu} — script \code{vm-power.sh}, journal d'exécution et explication des mesures de sécurité ajoutées.
\end{fiche}

\subsection{Le script \code{scripts/vm-power.sh}}
\begin{codebox}{scripts/vm-power.sh}
\begin{Verbatim}
#!/usr/bin/env bash
set -euo pipefail
RG="${1:-}"; ACTION="${2:-}"; ASSUME_YES="${3:-}"   # -y saute la confirmation

SCRIPT_DIR="$(cd "$(dirname "${BASH_SOURCE[0]}")" && pwd)"
LOG_FILE="$(dirname "$SCRIPT_DIR")/logs/vm-power.log"; mkdir -p "$(dirname "$LOG_FILE")"
log() { echo "$(date '+%Y-%m-%d %H:%M:%S') | $1" | tee -a "$LOG_FILE"; }

[[ -z "$RG" || -z "$ACTION" ]] && { echo "Usage: $0 <rg> <start|stop|deallocate|status> [-y]"; exit 1; }

# Garde-fou 1 : interdiction d'agir sur un RG de production
if [[ "$RG" == *prod* ]]; then
  log "REFUS : action '$ACTION' interdite sur un RG de production ('$RG')."; exit 1
fi
az group show --name "$RG" >/dev/null 2>&1 || { log "ERREUR : RG introuvable."; exit 1; }

# Garde-fou 2 : confirmation avant action destructive
confirm() {
  [[ "$ASSUME_YES" == "-y" ]] && return 0
  read -r -p "Confirmer '$ACTION' sur TOUTES les VM de '$RG' ? (oui/non) " ANS || true
  [[ "$ANS" == "oui" ]]
}
VMS=$(az vm list -g "$RG" --query "[].name" -o tsv)
case "$ACTION" in stop|deallocate)
  confirm || { log "Action '$ACTION' ANNULEE par l'utilisateur."; exit 0; } ;;
esac

log "=== Action '$ACTION' sur RG '$RG' ==="
for VM in $VMS; do
  case "$ACTION" in
    start)      az vm start      -g "$RG" -n "$VM" -o none; log "  start      -> $VM : OK" ;;
    stop)       az vm stop       -g "$RG" -n "$VM" -o none; log "  stop       -> $VM : OK" ;;
    deallocate) az vm deallocate -g "$RG" -n "$VM" -o none; log "  deallocate -> $VM : OK" ;;
    status)     STATE=$(az vm get-instance-view -g "$RG" -n "$VM" \
                  --query "instanceView.statuses[?starts_with(code,'PowerState/')].displayStatus" -o tsv)
                log "  status     -> $VM : $STATE" ;;
  esac
done
log "=== Termine ==="
\end{Verbatim}
\end{codebox}

\subsection{Mesures de sécurité ajoutées}
\begin{ltable}{Q[l,0.55] Q[l,1] X[l,1]}
Mesure & Mise en œuvre & Pourquoi \\
Garde-fou production & \code{if [[ "\$RG" == *prod* ]]} & Empêche tout arrêt accidentel d'un environnement de production. \\
Confirmation destructive & fonction \code{confirm()} sur \code{stop}/\code{deallocate} & Évite une coupure de service par erreur de frappe. \\
Bypass contrôlé & option \code{-y} & Permet l'automatisation (cron, CI) en connaissance de cause. \\
Journalisation & \code{log()} \arr{} \code{logs/vm-power.log} & Traçabilité : qui a fait quoi, quand. \\
Robustesse & \code{set -euo pipefail}, vérif du RG & Arrêt propre en cas d'erreur. \\
\end{ltable}

\subsection{Test de chaque action}
\begin{codebox}{console}
\begin{Verbatim}
$ ./vm-power.sh rg-shopeasy-prod deallocate
... | REFUS : action 'deallocate' interdite sur un RG de production. (code 1)

$ echo "non" | ./vm-power.sh rg-shopeasy-dev deallocate
... | Action 'deallocate' ANNULEE par l'utilisateur.

$ echo "oui" | ./vm-power.sh rg-shopeasy-dev stop
WARNING: About to power off the specified VM...
It will continue to be billed. To deallocate a VM, run: az vm deallocate.
... | stop       -> vm-shopeasy-dev-web-1 : OK
\end{Verbatim}
\end{codebox}
\begin{tipbox}[Confirmation FinOps par Azure lui-même]
Sur \code{stop}, Azure CLI affiche \textit{« It will continue to be billed. To deallocate a VM, run: az vm deallocate. »} — confirmant la distinction \code{stop} (facturé) / \code{deallocate} (compute libéré).
\end{tipbox}

\subsection{Journal d'exécution (\code{logs/vm-power.log})}
\begin{codebox}{logs/vm-power.log}
\begin{Verbatim}
2026-06-26 10:37:00 | status     -> vm-shopeasy-dev-web-1 : VM running
2026-06-26 10:37:03 | REFUS : action 'deallocate' interdite sur un RG de production.
2026-06-26 10:37:05 | Action 'deallocate' ANNULEE par l'utilisateur.
2026-06-26 10:38:23 | stop       -> vm-shopeasy-dev-web-1 : OK (VM arretee, toujours allouee)
2026-06-26 10:39:46 | deallocate -> vm-shopeasy-dev-web-1 : OK (compute libere)
2026-06-26 10:41:18 | status     -> vm-shopeasy-dev-web-1 : VM running
\end{Verbatim}
\end{codebox}
À l'issue des tests, l'environnement est remis à l'état nominal (les \textbf{2 VM \code{running}}).

\begin{etatbox}
\begin{itemize}
  \item Script \code{scripts/vm-power.sh} : 4 actions (\code{start}/\code{stop}/\code{deallocate}/\code{status}), toutes testées.
  \item 3 mesures de sécurité : garde-fou \code{prod}, confirmation destructive (+ bypass \code{-y}), journalisation horodatée.
  \item Distinction \code{stop} / \code{deallocate} confirmée par l'avertissement natif d'Azure CLI.
  \item Journal \code{logs/vm-power.log} produit (26 lignes) ; environnement remis à l'état nominal.
  \item \textbf{Prêt pour l'Atelier 7 — Exploiter un Storage Account.}
\end{itemize}
\end{etatbox}


% #####################################################################
%  ATELIER 7
% #####################################################################
\section{Atelier 7 — Exploiter un Storage Account}

\begin{fiche}
\textbf{\eTarget~Objectif} — disposer d'un espace de stockage opérationnel et sécurisé pour déposer les rapports et exports d'exploitation.

\textbf{\ePackage~Livrable attendu} — sortie CLI (et capture) montrant le conteneur, les fichiers déposés et l'accès public désactivé.
\end{fiche}

\subsection{Réutiliser et inspecter le Storage Account}
Le compte \code{shopeasydevdocs350wnq} (créé au TP2) est \textbf{réutilisé}. Inspection initiale :
\begin{codebox}{console}
\begin{Verbatim}
Name                   Sku           AllowBlobPublicAccess  MinimumTlsVersion  HttpsOnly
---------------------  ------------  ---------------------  -----------------  ---------
shopeasydevdocs350wnq  Standard_LRS  True                   TLS1_2             True
\end{Verbatim}
\end{codebox}
\textbf{Constat} : TLS 1.2 et HTTPS-only sont corrects, mais \textbf{\code{AllowBlobPublicAccess=True}} : le provider Terraform laisse ce paramètre activé par défaut — à corriger en exploitation.

\subsection{Sécuriser : désactiver l'accès public}
\begin{codebox}{bash}
\begin{Verbatim}
az storage account update --name "$STORAGE" -g "$RG" --allow-blob-public-access false
az storage container create --account-name "$STORAGE" --name "$CONTAINER" \
  --auth-mode login --public-access off
\end{Verbatim}
\end{codebox}
À partir d'ici, \textbf{aucun conteneur ne peut être rendu public}, même par erreur. Le conteneur privé \code{operations} est créé.

\subsection{RBAC : control-plane vs data-plane}
Le dépôt de blobs avec \code{--auth-mode login} a d'abord été \textbf{refusé} :
\begin{codebox}{console}
\begin{Verbatim}
ERROR: You do not have the required permissions needed to perform this operation.
    "Storage Blob Data Contributor" ...
\end{Verbatim}
\end{codebox}
Être \textit{Owner} de l'abonnement donne les droits sur le \textbf{control-plane} (créer/configurer via ARM), mais \textbf{pas} sur le \textbf{data-plane} (lire/écrire le contenu des blobs). La bonne pratique — conforme au moindre privilège — est d'\textbf{assigner un rôle RBAC data-plane} précis, au scope du compte (pas de clé partagée) :
\begin{codebox}{bash}
\begin{Verbatim}
az role assignment create --role "Storage Blob Data Contributor" \
  --assignee-object-id "$USER_OID" --assignee-principal-type User \
  --scope "$STORAGE_ID"
\end{Verbatim}
\end{codebox}

\subsection{Déposer les exports et vérifier la fermeture}
\begin{codebox}{console}
\begin{Verbatim}
$ az storage blob list --account-name "$STORAGE" -c operations --auth-mode login -o table
Nom                                         TailleOctets
------------------------------------------  ------------
inventaires/resources.json                  12014
inventaires/resources.tsv                   983
inventaires/vms-20260626-103147.txt         281
rapports/rapport-exploitation-shopeasy.md   9772

$ curl -s -o /dev/null -w "%{http_code}" \
    "https://$STORAGE.blob.core.windows.net/operations/inventaires/resources.tsv"
409   ->   <Code>PublicAccessNotPermitted</Code>
\end{Verbatim}
\end{codebox}
\begin{notebox}[Preuve d'accès fermé]
Même en connaissant l'\textbf{URL exacte}, un visiteur anonyme reçoit \textbf{409 \code{PublicAccessNotPermitted}} : le stockage est bien fermé. L'accès n'est possible qu'avec une \textbf{identité Azure autorisée} (RBAC).
\end{notebox}

\subsection{Pourquoi le stockage de rapports ne doit pas être public}
\begin{itemize}
  \item \textbf{Cartographie offerte à un attaquant} : un rapport liste noms de ressources, IP, état des VM, règles réseau.
  \item \textbf{Fuite de données / RGPD} : un blob public est accessible à quiconque connaît l'URL, sans traçabilité.
  \item \textbf{Moindre privilège} : seuls les exploitants \textbf{authentifiés} (RBAC) doivent lire ces rapports.
  \item \textbf{Conformité} : « pas de stockage public » est un contrôle standard (CIS Azure).
\end{itemize}

\subsection{Captures}
\screenshot{atelier07-storage-public-access.png}{Configuration du compte (Portal \arr{} Storage accounts \arr{} \dots{} \arr{} Configuration) : \textit{Allow blob anonymous access} = \textbf{Disabled}.}
\screenshot{atelier07-container-operations.png}{Conteneur \code{operations} et le dossier \code{inventaires/} (exports déposés). Le niveau « Private » n'apparaît plus car l'accès anonyme est désactivé au niveau du compte.}

\begin{etatbox}
\begin{itemize}
  \item Storage Account réutilisé et \textbf{sécurisé} : accès public \textbf{désactivé} (était \code{True}), TLS 1.2, HTTPS-only.
  \item Conteneur privé \code{operations} créé ; 4 exports déposés sous \code{inventaires/}.
  \item RBAC data-plane appliqué proprement (\code{Storage Blob Data Contributor}, scope = compte) — pas de clé partagée.
  \item Accès anonyme \textbf{prouvé refusé} (HTTP 409 \code{PublicAccessNotPermitted}).
  \item \textbf{Prêt pour l'Atelier 8 — Surveiller les métriques avec Azure Monitor.}
\end{itemize}
\end{etatbox}


% #####################################################################
%  ATELIER 8
% #####################################################################
\section{Atelier 8 — Surveiller les métriques avec Azure Monitor}

\begin{fiche}
\textbf{\eTarget~Objectif} — interroger les métriques Azure et créer une alerte simple sur une ressource critique.

\textbf{\ePackage~Livrable attendu} — sortie CLI montrant la métrique CPU et la règle d'alerte créée, avec justification du seuil.
\end{fiche}

\subsection{Identifier la VM et lister les métriques}
\begin{codebox}{bash}
\begin{Verbatim}
VM_ID=$(az vm show -g "$RG" -n "$VM1" --query id -o tsv)
az monitor metrics list-definitions --resource "$VM_ID" \
  --query "[].{Metrique:name.value, Unite:unit}" -o table
\end{Verbatim}
\end{codebox}
\begin{codebox}{console}
\begin{Verbatim}
Metrique                     Unite
---------------------------  --------------
Percentage CPU               Percent
Network In / Network Out     Bytes
Disk Read Bytes / Write      Bytes
CPU Credits Remaining        Count
CPU Credits Consumed         Count
...                          (~30 metriques au total)
\end{Verbatim}
\end{codebox}
\begin{tipbox}[Spécificité Standard\_B2ats\_v2 (burstable)]
C'est une VM \textit{burstable} (famille B). Les métriques \code{CPU Credits Remaining} et \code{CPU Credits Consumed} sont propres à cette famille : si le CPU reste élevé, la VM épuise ses crédits et se fait \textbf{brider} (throttling).
\end{tipbox}

\subsection{Lire la métrique CPU}
\begin{codebox}{console}
\begin{Verbatim}
Heure                 CPU_moyen_pct
--------------------  ---------------
2026-06-26T08:42:00Z  0.15
2026-06-26T08:46:00Z  0.62
2026-06-26T08:55:00Z  0.12
2026-06-26T09:00:00Z  0.15
\end{Verbatim}
\end{codebox}
Le CPU moyen oscille entre \textbf{0,12\,\% et 0,62\,\%} : la VM web est \textbf{au repos} (pas de trafic). La métrique fournit une \textbf{base de référence} pour calibrer le seuil d'alerte.

\subsection{Créer une alerte CPU > 80\,\%}
\begin{codebox}{bash}
\begin{Verbatim}
az monitor metrics alert create --name "alert-cpu-high-$VM1" -g "$RG" \
  --scopes "$VM_ID" --condition "avg Percentage CPU > 80" \
  --description "Alerte CPU elevee sur VM web ShopEasy" \
  --evaluation-frequency 1m --window-size 5m --severity 3
\end{Verbatim}
\end{codebox}
\begin{codebox}{console}
\begin{Verbatim}
Nom                                   Active    Severite    Condition
------------------------------------  --------  ----------  -----------
alert-cpu-high-vm-shopeasy-dev-web-1  True      3           GreaterThan
\end{Verbatim}
\end{codebox}

\subsection{Justification du seuil}
\begin{ltable}{Q[l,0.4] Q[l,0.5] X[l,1]}
Paramètre & Valeur & Justification \\
Seuil & \code{> 80 \%} & Assez haut pour ignorer les pics normaux, assez bas pour détecter une saturation avant dégradation. \\
Agrégation & \code{Average} & On surveille une charge \textbf{soutenue}, pas un pic instantané. \\
Fenêtre & \code{5 min} & \textbf{Lisse} les pics transitoires (charge élevée prolongée requise). \\
Fréquence & \code{1 min} & Réactivité : réévaluation chaque minute. \\
Sévérité & \code{3} (Low) & CPU élevé = signal d'attention, pas une panne (Sev 0/1 = indispo). \\
\end{ltable}

\subsection{Deux autres alertes utiles + limites}
\begin{ltable}{Q[l,0.6] X[l,1] Q[c,0.4]}
Alerte proposée & Pourquoi & Sévérité \\
Disponibilité backends (LB \code{Health Probe Status}) & Une VM ne répond plus à la sonde HTTP : indicateur proche de l'expérience client. & 1 \\
Saturation disque / mémoire (VM) & Disque plein (logs Nginx) ou mémoire saturée bloque le service. & 2 \\
\end{ltable}
\begin{notebox}[Limites d'une alerte mal calibrée]
\textbf{Trop sensible} (seuil bas, fenêtre courte) \arr{} faux positifs, \textit{fatigue d'alerte}, vrais incidents ratés. \textbf{Trop large} (seuil haut, fenêtre longue) \arr{} détection tardive. Une alerte doit être \textbf{actionnable} et associée à une \textbf{procédure} (runbook).
\end{notebox}

\subsection{Captures}
\screenshot{atelier08-alerte-cpu.png}{Règle d'alerte (Portal \arr{} Monitor \arr{} Alerts \arr{} Alert rules) : \code{alert-cpu-high-vm-shopeasy-dev-web-1}, \textit{Enabled}, sévérité 3, condition \textit{Percentage CPU > 80}.}
\screenshot{atelier08-metrique-cpu.png}{Graphe de la métrique \textit{Percentage CPU} (Portal \arr{} Virtual machines \arr{} \dots{} \arr{} Metrics). Le graphe reste très bas (VM au repos).}

\begin{etatbox}
\begin{itemize}
  \item Métriques de la VM identifiées (CPU, réseau, disque, crédits CPU burstable) ; CPU lu (baseline 0,12–0,62\,\%).
  \item Alerte \code{alert-cpu-high-vm-shopeasy-dev-web-1} créée et active : \code{avg Percentage CPU > 80}, fenêtre 5 min, fréquence 1 min, sévérité 3.
  \item Seuil justifié ; deux alertes complémentaires proposées (disponibilité LB, saturation disque/mémoire).
  \item Limites des alertes mal calibrées documentées (fatigue d'alerte vs détection tardive).
  \item \textbf{Prêt pour l'Atelier 9 — Écrire un script de contrôle de santé.}
\end{itemize}
\end{etatbox}


% #####################################################################
%  ATELIER 9
% #####################################################################
\section{Atelier 9 — Écrire un script de contrôle de santé}

\begin{fiche}
\textbf{\eTarget~Objectif} — créer un script qui vérifie rapidement si l'environnement est dans un état acceptable.

\textbf{\ePackage~Livrable attendu} — script \code{healthcheck.sh} final, preuve d'exécution et commentaires sur les vérifications ajoutées.
\end{fiche}

\subsection{Le script \code{scripts/healthcheck.sh}}
Huit contrôles (4 de base + 5 améliorations), code de sortie \code{0} (conforme) ou \code{1} (avertissement).
\begin{codebox}{scripts/healthcheck.sh}
\begin{Verbatim}
#!/usr/bin/env bash
set -euo pipefail
RG="${1:-rg-shopeasy-dev}"; CONTAINER="${2:-operations}"; STATUS=0
ok()   { echo "  [OK]   $1"; }
warn() { echo "  [WARN] $1"; STATUS=1; }
ko()   { echo "  [KO]   $1"; STATUS=1; }

# 1. Groupe de ressources (bloquant)
az group show --name "$RG" >/dev/null 2>&1 && ok "groupe '$RG' present" \
  || { ko "RG introuvable"; exit 1; }
# 2. Au moins une VM
NB_VM=$(az vm list -g "$RG" --query "length(@)" -o tsv)
[[ "$NB_VM" -ge 1 ]] && ok "$NB_VM VM(s) presente(s)" || warn "aucune VM"
# 3. Tag Owner sur le groupe
OWNER=$(az group show --name "$RG" --query "tags.Owner" -o tsv)
[[ -n "$OWNER" ]] && ok "tag Owner present ($OWNER)" || warn "tag Owner absent"
# 4. Ressources sans tag Application
UNTAGGED=$(az resource list -g "$RG" --query "length([?tags.Application==null])" -o tsv)
[[ "$UNTAGGED" -eq 0 ]] && ok "toutes taguees Application" || warn "$UNTAGGED sans tag Application"
# 5. Alertes Azure Monitor
ALERTS=$(az monitor metrics alert list -g "$RG" --query "length(@)" -o tsv)
[[ "$ALERTS" -ge 1 ]] && ok "$ALERTS alerte(s)" || warn "aucune alerte configuree"
# 6. Storage Account (recupere dynamiquement)
STORAGE=$(az storage account list -g "$RG" --query "[0].name" -o tsv 2>/dev/null || true)
[[ -n "$STORAGE" ]] && ok "Storage present ($STORAGE)" || warn "aucun Storage Account"
# 7. Conteneur operations
az storage container show --account-name "$STORAGE" --name "$CONTAINER" \
  --auth-mode login >/dev/null 2>&1 && ok "conteneur '$CONTAINER' present" \
  || warn "conteneur '$CONTAINER' absent"
# 8. Regle NSG HTTP/HTTPS
NB_HTTP=$(az network nsg list -g "$RG" \
  --query "length([].securityRules[?direction=='Inbound' && access=='Allow' && \
           (destinationPortRange=='80' || destinationPortRange=='443')][])" -o tsv)
[[ "$NB_HTTP" -ge 1 ]] && ok "$NB_HTTP regle(s) NSG HTTP/HTTPS" || warn "aucune regle HTTP/HTTPS"

exit $STATUS
\end{Verbatim}
\end{codebox}

\subsection{Preuve d'exécution}
\begin{codebox}{console}
\begin{Verbatim}
1. Groupe de ressources    [OK]   groupe 'rg-shopeasy-dev' present
2. Machines virtuelles      [OK]   2 VM(s) presente(s)
3. Gouvernance - Owner      [OK]   tag Owner present (cloudops-shopeasy)
4. Gouvernance - Application[WARN] 11 ressource(s) sans tag Application
5. Supervision - alertes    [OK]   1 alerte(s) configuree(s)
6. Stockage - compte        [OK]   Storage present (shopeasydevdocs350wnq)
7. Stockage - operations    [OK]   conteneur 'operations' present
8. Reseau - NSG HTTP/HTTPS  [OK]   1 regle(s) NSG autorisant HTTP/HTTPS
 RESULTAT : ATTENTION - points a verifier (voir [WARN])     >>> Code de sortie : 1
\end{Verbatim}
\end{codebox}
\textbf{7 contrôles sur 8 au vert.} Le seul avertissement — \textbf{11 ressources sans tag \code{Application}} — est cohérent avec l'Atelier 3. Le code \textbf{1} permet d'intégrer le script en \textbf{CI/cron} (un code $\neq$ 0 déclenche une notification).

\subsection{Vérifications ajoutées}
\begin{ltable}{Q[l,0.55] X[l,1]}
Contrôle ajouté & Pourquoi \\
Au moins une VM & Un groupe sans VM = service web absent. \\
Tag \code{Owner} & Sans responsable identifié, pas de gouvernance (qui contacter ?). \\
Storage Account & Sans stockage, impossible d'archiver les rapports. \\
Conteneur \code{operations} & Valide l'espace de dépôt des rapports. \\
Règle NSG HTTP/HTTPS & Sans ouverture 80/443, le site est injoignable. \\
\end{ltable}

\begin{etatbox}
\begin{itemize}
  \item Script \code{scripts/healthcheck.sh} final : \textbf{8 contrôles}, code de sortie 0/1.
  \item Exécution réelle : \textbf{7 OK / 1 WARN} (11 ressources sans tag \code{Application}), code 1 \arr{} intégrable en CI/cron.
  \item Contrôles de gouvernance, stockage et réseau ajoutés et commentés ; Storage récupéré dynamiquement.
  \item \textbf{Prêt pour l'Atelier 10 — Automatiser avec Python et Azure SDK.}
\end{itemize}
\end{etatbox}


% #####################################################################
%  ATELIER 10
% #####################################################################
\section{Atelier 10 — Automatiser avec Python et Azure SDK}

\begin{fiche}
\textbf{\eTarget~Objectif} — découvrir le SDK Azure pour Python afin de produire un inventaire programmable.

\textbf{\ePackage~Livrable attendu} — script \code{inventory.py} et fichier CSV généré, lisible dans Excel ou LibreOffice Calc.
\end{fiche}

\subsection{Installation des dépendances}
\begin{codebox}{bash}
\begin{Verbatim}
python3 -m venv .venv && source .venv/bin/activate
pip install azure-identity azure-mgmt-resource azure-mgmt-compute
\end{Verbatim}
\end{codebox}
\begin{codebox}{console}
\begin{Verbatim}
azure-identity 1.25.3 | azure-mgmt-resource 26.0.0 | azure-mgmt-compute 38.1.0
\end{Verbatim}
\end{codebox}

\subsection{Le script \code{python/inventory.py}}
\begin{codebox}{python/inventory.py}
\begin{Verbatim}
#!/usr/bin/env python3
"""inventory.py - Inventaire programmable de ShopEasy via l'Azure SDK -> CSV."""
import os, csv
from azure.identity import DefaultAzureCredential
from azure.mgmt.compute import ComputeManagementClient
try:
    from azure.mgmt.resource import ResourceManagementClient            # SDK < 26
except ImportError:
    from azure.mgmt.resource.resources import ResourceManagementClient  # SDK >= 26

subscription_id = os.environ.get("AZURE_SUBSCRIPTION_ID")
resource_group = os.environ.get("AZURE_RESOURCE_GROUP", "rg-shopeasy-dev")
if not subscription_id:
    raise SystemExit("Variable AZURE_SUBSCRIPTION_ID manquante")

credential = DefaultAzureCredential()
resource_client = ResourceManagementClient(credential, subscription_id)
compute_client = ComputeManagementClient(credential, subscription_id)

def deduce_role(name, res_type):
    t = res_type.lower(); n = name.lower()
    if "virtualmachines" in t:   return "Serveur web (Ubuntu + Nginx)"
    if "loadbalancers" in t:     return "Repartiteur de charge (point d'entree HTTP)"
    if "publicipaddresses" in t: return "IP publique du LB" if "lb" in n else "IP publique VM"
    if "storageaccounts" in t:   return "Stockage documentaire prive"
    if "insights" in t:          return "Alerte de supervision (Azure Monitor)"
    return "Autre"   # (+ VNet, NSG, NIC, disks)

def fmt_tags(tags):
    return ";".join(f"{k}={v}" for k, v in sorted(tags.items())) if tags else "(aucun)"

vm_sizes = {vm.name.lower(): vm.hardware_profile.vm_size
            for vm in compute_client.virtual_machines.list(resource_group)}
rows = []
for res in resource_client.resources.list_by_resource_group(resource_group):
    rows.append({"Nom": res.name, "Type": res.type, "Region": res.location,
                 "Taille": vm_sizes.get(res.name.lower(), ""),
                 "Tags": fmt_tags(res.tags), "Role": deduce_role(res.name, res.type)})
rows.sort(key=lambda r: (r["Type"], r["Nom"]))
with open("exports/inventory.csv", "w", newline="", encoding="utf-8-sig") as f:
    w = csv.DictWriter(f, fieldnames=["Nom","Type","Region","Taille","Tags","Role"])
    w.writeheader(); w.writerows(rows)
print(f"Ressources inventoriees: {len(rows)} | VM: {len(vm_sizes)} | CSV genere")
\end{Verbatim}
\end{codebox}

\begin{notebox}[Adaptation au SDK récent]
Le code du sujet (\code{from azure.mgmt.resource import ResourceManagementClient}) \textbf{échoue} avec \code{azure-mgmt-resource 26.0.0} : le client a été déplacé dans le sous-module \code{azure.mgmt.resource.resources}. Le script gère les deux cas via un \code{try/except ImportError} (rétrocompatible).
\end{notebox}

\subsection{Exécution et CSV généré}
\code{DefaultAzureCredential} réutilise la session \code{az login} (via \code{AzureCliCredential}), sans secret.
\begin{codebox}{console}
\begin{Verbatim}
Inventaire du groupe   : rg-shopeasy-dev
Ressources inventoriees: 14
Machines virtuelles    : 2
CSV genere             : exports/inventory.csv
\end{Verbatim}
\end{codebox}
\begin{codebox}{exports/inventory.csv}
\begin{Verbatim}
Nom,Type,Region,Taille,Tags,Role
vm-shopeasy-dev-web-1,Microsoft.Compute/virtualMachines,swedencentral,Standard_B2ats_v2,
  Application=shopeasy;Environment=dev;Owner=cloudops-shopeasy;Role=web,Serveur web (Ubuntu + Nginx)
alert-cpu-high-vm-shopeasy-dev-web-1,Microsoft.Insights/metricalerts,global,,(aucun),
  Alerte de supervision (Azure Monitor)
lb-shopeasy-dev-web,Microsoft.Network/loadBalancers,swedencentral,,...,Repartiteur de charge
\end{Verbatim}
\end{codebox}
\begin{tipbox}[14 ressources via le SDK, 13 via la CLI]
Le SDK Python retourne l'\textbf{alerte métrique} \code{Microsoft.Insights/metricAlerts} (créée à l'Atelier 8), que \code{az resource list} \textbf{n'inclut pas} dans son décompte. Ce n'est pas une incohérence mais une \textbf{différence de comportement} documentée entre les deux outils.
\end{tipbox}

\subsection{Bash ou Python ?}
\begin{ltable}{Q[l,1] X[l,1]}
Besoin & Outil adapté \\
Enchaîner des commandes, contrôles rapides, exports bruts & \textbf{Bash} (\code{inventory.sh}, \code{healthcheck.sh}) \\
Enrichir / croiser des données, format structuré (CSV), intégrer une API & \textbf{Python + SDK} (\code{inventory.py}) \\
\end{ltable}
Ici, le \textbf{croisement} ressources $\leftrightarrow$ tailles de VM et la \textbf{déduction de rôle} illustrent le moment où Python devient préférable : le besoin n'est plus « lister » mais « \textbf{structurer et enrichir} ».

\subsection{Capture}
\screenshot{atelier10-csv-tableur.png}{Le fichier \code{exports/inventory.csv} ouvert dans un tableur : les colonnes \textit{Nom, Type, Region, Taille, Tags, Role} s'affichent correctement (encodage \code{utf-8-sig}).}

\begin{etatbox}
\begin{itemize}
  \item Environnement Python isolé (\code{.venv}) ; SDK Azure installé (identity, mgmt-resource, mgmt-compute).
  \item Script \code{python/inventory.py} fonctionnel : authentification \code{az login}, \textbf{CSV enrichi} (rôle + taille).
  \item Adaptation au SDK récent gérée (\code{ResourceManagementClient} déplacé en v26, import rétrocompatible).
  \item \code{exports/inventory.csv} (14 ressources) lisible dans un tableur ; classe l'alerte de supervision.
  \item \textbf{Prêt pour l'Atelier 11 — Analyse FinOps d'exploitation.}
\end{itemize}
\end{etatbox}


% #####################################################################
%  ATELIER 11
% #####################################################################
\section{Atelier 11 — Analyse FinOps d'exploitation}

\begin{fiche}
\textbf{\eTarget~Objectif} — identifier les actions simples permettant de limiter les coûts d'un environnement de développement.

\textbf{\ePackage~Livrable attendu} — tableau FinOps complété et trois recommandations courtes pour la DSI.
\end{fiche}

\subsection{Ressources les plus coûteuses}
Coûts calculés à partir des \textbf{prix réels} de l'API Azure Retail Prices (\code{swedencentral}, USD, base 730\,h/mois).
\begin{ltable}{Q[l,0.7] Q[l,0.7] Q[c,0.45] Q[c,0.35]}
Poste & Détail & Coût/mois & Part \\
\textbf{Load Balancer Standard} & forfait règles + data & $\approx$ 18,25 \$ & $\sim$39 \% \\
\textbf{Compute — 2 VM \code{B2ats\_v2}} & 2 $\times$ 0,00972 \$/h $\times$ 730 & $\approx$ 14,19 \$ & $\sim$30 \% \\
\textbf{3 IP publiques Standard} & 3 $\times$ 0,005 \$/h $\times$ 730 & $\approx$ 10,95 \$ & $\sim$23 \% \\
2 disques OS \code{Standard\_LRS} & $\sim$30 Go HDD & $\approx$ 2,60 \$ & $\sim$6 \% \\
Storage Account LRS & usage faible & $<$ 1 \$ & $\sim$2 \% \\
\textbf{Total} & & \textbf{$\approx$ 47 \$} & 100 \% \\
\end{ltable}
Les trois premiers postes représentent \textbf{$\sim$92 \%} du coût. Le \textbf{Load Balancer} et les \textbf{IP} sont facturés \textbf{24h/24, indépendamment du trafic} ; seul le \textbf{compute} est élastique.
\begin{tipbox}[Aucune ressource orpheline]
\code{az disk list --query "[?managedBy==null]"} et \code{az network public-ip list --query "[?ipConfiguration==null]"} \arr{} \textbf{aucun disque orphelin, aucune IP non associée}. L'environnement est propre, sans coût « invisible ».
\end{tipbox}

\subsection{Tableau FinOps}
\begin{ltable}{Q[l,0.6] Q[l,0.7] Q[l,0.7] X[l,1]}
Action & Effet & Risque / limite & Application à ShopEasy \\
Désallouer les VM hors usage & Réduction du coût compute & Indisponibilité & \code{vm-power.sh} : VM2 désallouée par défaut, VM1 hors heures. \\
Réduire la taille des VM & Réduction du coût mensuel & Performance plus faible & \code{B2ats\_v2} déjà petite ; surveiller les crédits CPU. \\
Supprimer les disques inutilisés & Réduction du stockage & Perte de donnée utile & Aucun disque orphelin ; contrôle automatisé. \\
Standardiser les tags & Meilleure analyse des coûts & Discipline nécessaire & Étendre aux 11 ressources restantes. \\
Définir un budget & Alerte en cas de dérive & Ne bloque pas la dépense & Budget 50 \$/mois, alertes 80 / 100 \%. \\
\end{ltable}

\subsection{Stratégie d'arrêt/démarrage \& budget}
\begin{itemize}
  \item \textbf{Désallocation hors heures ouvrées} (\code{vm-power.sh deallocate} le soir, \code{start} le matin) : une VM utilisée $\sim$50\,h/semaine économise \textbf{$\sim$70 \%} de son compute.
  \item \textbf{Automatisation} : Azure Automation runbook, Logic App ou cron ; tag \code{AutoShutdown}.
  \item \textbf{VM secondaire} \code{web-2} (HA) : \textbf{désallouée par défaut} en dev.
  \item \textbf{Levier maximal} : \code{terraform destroy} entre les séances \arr{} coût $\approx$ 0, recréable par \code{apply}.
  \item \textbf{Budget} : 50 \$/mois sur le RG, alertes \textbf{80 \% (40 \$)} et \textbf{100 \% (50 \$)} — un budget alerte mais ne bloque pas.
\end{itemize}

\subsection{Trois recommandations FinOps pour la DSI}
\begin{ltable}{Q[c,0.2] X[l,1.3] Q[l,0.7]}
\# & Recommandation & Gain \\
1 & \textbf{Désallocation automatisée} des VM dev + destruction entre séances & $-60$ à $-70$ \% compute \\
2 & \textbf{Supprimer les IP publiques des VM}, administrer via \textbf{Azure Bastion} & $\approx$ 7 \$/mois + sécurité \\
3 & \textbf{Imposer les tags par Azure Policy} + budget mensuel avec alertes & gouvernance, détection des dérives \\
\end{ltable}

\begin{etatbox}
\begin{itemize}
  \item Coût réel chiffré (prix Azure Retail Prices) : \textbf{$\approx$ 47 \$/mois en 24/7}, dont $\sim$92 \% sur LB + compute + IP.
  \item Environnement \textbf{propre} : aucun disque orphelin, aucune IP non associée.
  \item Tableau FinOps complété (5 actions $\times$ effet $\times$ risque $\times$ application ShopEasy).
  \item Stratégie d'arrêt, politique de tags, budget (50 \$) et \textbf{3 recommandations DSI} prêtes.
  \item \textbf{Prêt pour l'Atelier 12 — Rapport d'exploitation ShopEasy.}
\end{itemize}
\end{etatbox}


% #####################################################################
%  ATELIER 12 — RAPPORT D'EXPLOITATION
% #####################################################################
\section{Atelier 12 — Rapport d'exploitation ShopEasy}

\begin{tcolorbox}[enhanced, sharp corners, boxrule=0pt, leftrule=4pt, colback=azurepale, colframe=azuredark,
   left=4.5mm,right=4.5mm,top=3.6mm,bottom=3.6mm, before skip=10pt, after skip=12pt]
\textbf{Objet} — exploitation et automatisation de l'environnement Azure de ShopEasy\\[2pt]
\textbf{Périmètre} — groupe de ressources \code{rg-shopeasy-dev} (\code{swedencentral}), abonnement Azure for Students\\[2pt]
\textbf{Auteurs (binôme)} — Louis \textsc{Scarfone} \& Maxence \textsc{Bourragué}\\[2pt]
\textbf{Double lecture} — l'équipe technique (commandes rejouables) et la DSI (risques, coûts, priorités)
\end{tcolorbox}

\subsection{Contexte et périmètre}
ShopEasy a conçu son architecture au TP1 (déploiement manuel) puis l'a industrialisée au TP2 (Terraform). Le TP3 ouvre la phase d'\textbf{exploitation} : administrer, superviser, sécuriser et optimiser avec \textbf{Azure CLI, Bash et Python}. L'environnement a été \textbf{redéployé à l'identique} depuis le code Terraform du TP2 (21 objets), puis exploité sans repasser par le portail.

\subsection{Inventaire des ressources}
Inventaire produit en CLI, exporté en JSON/TSV, enrichi en CSV par le SDK Python. \textbf{13 ressources d'infrastructure} (décompte \code{az resource list}) + \textbf{1 alerte de supervision} retournée par le SDK, soit \textbf{14 objets ARM} : 2 VM web, 2 disques, 3 IP publiques, 2 NIC, 1 VNet, 1 NSG, 1 Load Balancer, 1 Storage Account, 1 alerte métrique. \textbf{Statut : conforme.}

\subsection{Tags et gouvernance}
Convention d'exploitation \textbf{normalisée} (\code{Application}, \code{Environment}, \code{Owner}, \code{CostCenter}, \code{ManagedBy}) appliquée au RG et aux 2 VM, en remplacement des tags techniques Terraform. \textbf{11 ressources} conservent encore les tags \code{snake\_case} — écart détecté automatiquement par \code{inventory.sh} et \code{healthcheck.sh}. \textbf{Statut : partiel.}

\subsection{État des VM et actions réalisées}
Administration complète en CLI, sans portail ni SSH direct : \code{run-command} (Nginx \code{active}, HTTP 200), \code{restart} sur \code{web-1}, \code{deallocate} puis \code{start} sur \code{web-2}. Distinction \code{stop} / \code{deallocate} maîtrisée. \textbf{Statut : conforme.}

\subsection{Stockage opérationnel}
Storage Account du TP2 \textbf{réutilisé et sécurisé} : accès public \textbf{désactivé} (\code{AllowBlobPublicAccess=False}), TLS 1.2, HTTPS-only ; conteneur \textbf{privé \code{operations}} ; authentification par \textbf{RBAC data-plane} (sans clé partagée) ; accès anonyme \textbf{prouvé refusé} (HTTP 409). \textbf{Statut : conforme.}

\subsection{Monitoring et alertes}
Métriques identifiées (CPU, réseau, disque, crédits CPU) ; \textbf{alerte active} \code{alert-cpu-high-vm-shopeasy-dev-web-1} (\code{avg Percentage CPU > 80}, fenêtre 5 min, sévérité 3) ; deux alertes complémentaires \textbf{proposées} (disponibilité LB, saturation disque). \textbf{Statut : partiel.}

\subsection{Scripts produits}
\begin{ltable}{Q[l,0.55] X[l,1]}
Script & Rôle \\
\code{inventory.sh} & Inventaire + synthèse (compteurs, ressources sans tag, avertissement VM \code{running}). \\
\code{vm-power.sh} & Start/stop/deallocate/status, garde-fou \code{prod}, confirmation, journal. \\
\code{healthcheck.sh} & 8 contrôles de santé, code de sortie 0/1 (CI/cron). \\
\code{inventory.py} & Inventaire SDK $\rightarrow$ CSV (rôle déduit, taille VM, \code{utf-8-sig}). \\
\code{variables.sh} & Variables centralisées, aucun secret en dur. \\
\end{ltable}
Tous commentés, paramétrables, relançables, sans action destructive non confirmée. \textbf{Statut : conforme.}

\subsection{Analyse sécurité}
\textbf{Contrôles en place} : SSH restreint à l'IP admin (\code{/32}) ; HTTP limité au port 80 ; stockage privé + RBAC (moindre privilège) ; aucun secret versionné ; administration sans SSH direct (\code{run-command}). \textbf{Risques résiduels} : IP publiques directes sur les VM, absence d'Azure Bastion, journalisation/diagnostic non centralisée. \textbf{Statut : conforme} sur l'essentiel, renforcements identifiés.

\subsection{Analyse FinOps}
Coût réel \textbf{$\approx$ 47 \$/mois en 24/7} : Load Balancer ($\sim$39 \%), compute ($\sim$30 \%), IP publiques ($\sim$23 \%). Le LB et les IP facturent même à trafic nul ; seul le compute est élastique. Aucune ressource orpheline. \textbf{Statut : partiel} — coût chiffré et leviers identifiés, budget et désallocation auto à appliquer.

\subsection{Recommandations d'amélioration (priorisées)}
\begin{ltable}{Q[c,0.2] X[l,1.3] Q[l,0.55] Q[c,0.4]}
\# & Recommandation & Gain & Priorité \\
1 & Désallocation automatisée + destruction entre séances & $-60$ à $-70$ \% compute & Haute \\
2 & Supprimer les IP publiques des VM (Azure Bastion) & $\approx$ 7 \$/mois + sécurité & Haute \\
3 & Généraliser les tags + Azure Policy & gouvernance & Moyenne \\
4 & Budget (50 \$/mois) avec alertes 80/100 \% & détection des dérives & Moyenne \\
5 & Compléter les alertes (LB, disque) + action group & détection des pannes & Moyenne \\
6 & Diagnostic settings / Activity Log centralisés & auditabilité & Basse \\
\end{ltable}

\subsection*{Tableau de synthèse}
\begin{ltable}{Q[l,0.5] Q[l,0.5] X[l,1.3]}
Domaine & Statut & Commentaire \\
\textbf{Inventaire} & \eGreen~Conforme & Inventaire CLI complet (14 objets), exports JSON/TSV/CSV, rejouable. \\
\textbf{Tags} & \eYellow~Partiel & RG + 2 VM normalisés ; 11 ressources encore en \code{snake\_case}. \\
\textbf{VM} & \eGreen~Conforme & États, \code{run-command}, \code{start}/\code{stop}/\code{deallocate} maîtrisés. \\
\textbf{Stockage} & \eGreen~Conforme & Conteneur privé, accès public désactivé, RBAC, accès anonyme refusé (409). \\
\textbf{Monitoring} & \eYellow~Partiel & Alerte CPU active ; alertes disponibilité/disque et action group à créer. \\
\textbf{Sécurité} & \eGreen~Conforme & SSH restreint, stockage privé, RBAC, aucun secret versionné ; Bastion recommandé. \\
\textbf{FinOps} & \eYellow~Partiel & Coût chiffré ($\sim$47 \$/mois), leviers identifiés ; budget et désallocation à appliquer. \\
\end{ltable}

\subsection*{Conclusion}
L'environnement ShopEasy est désormais \textbf{exploité de manière outillée et reproductible} : inventaire automatisé, administration scriptée et sécurisée, stockage verrouillé, supervision initiale et analyse de coûts chiffrée. Les fondamentaux (inventaire, VM, stockage, sécurité) sont \textbf{conformes} ; les axes \textbf{tags}, \textbf{monitoring} et \textbf{FinOps} sont engagés et \textbf{partiellement} appliqués, avec une feuille de route claire. L'ensemble constitue un \textbf{mini-runbook} permettant à une autre personne de reprendre l'administration \textbf{sans dépendre du portail Azure}.

\begin{etatbox}[Fin du TP3]
\begin{itemize}
  \item Rapport d'exploitation complet (10 sections + tableau de synthèse), déposé dans le conteneur \code{operations/rapports/}.
  \item Bilan : \textbf{4 domaines conformes} (inventaire, VM, stockage, sécurité), \textbf{3 partiels} (tags, monitoring, FinOps) avec recommandations.
  \item Mini-kit d'exploitation livré : \code{variables.sh}, 3 scripts Bash, \code{inventory.py}, exports, logs et rapport.
\end{itemize}
\end{etatbox}


% #####################################################################
%  QUIZ
% #####################################################################
\section{Quiz de validation}

\begin{enumerate}[leftmargin=2.2em]
\item \textbf{Différence entre \code{az vm stop} et \code{az vm deallocate} ?}\\ \code{stop} arrête l'OS mais la VM reste \textbf{allouée} (compute facturé) ; \code{deallocate} arrête \textbf{et libère} le calcul (\textit{Stopped (deallocated)}, compute non facturé). En dev, privilégier \code{deallocate}.
\item \textbf{Pourquoi Azure CLI plutôt que le portail pour une tâche répétitive ?}\\ Une commande est \textbf{rejouable}, documentable, intégrable à un script et exécutable en pipeline. Le portail est manuel, peu reproductible et difficile à tracer.
\item \textbf{À quoi sert l'option \code{--query} ?}\\ À filtrer et transformer les sorties JSON avec le langage \textbf{JMESPath} (sélectionner des champs, filtrer, renommer des colonnes).
\item \textbf{Quel format de sortie pour réutiliser un résultat dans un script Bash ?}\\ Le \textbf{TSV} (\code{-o tsv}) : sans guillemets ni en-tête, il s'extrait directement dans une variable Bash.
\item \textbf{Pourquoi les tags sont-ils importants pour le FinOps ?}\\ Ils permettent de \textbf{ventiler la facture} (application, environnement, \code{CostCenter}), d'identifier le responsable (\code{Owner}) et d'automatiser des politiques (arrêt des \code{dev}).
\item \textbf{Qu'est-ce qu'un runbook d'exploitation ?}\\ Une \textbf{procédure documentée} décrivant comment réaliser une action de façon fiable (démarrer/arrêter, vérifier la santé, réagir à une alerte). Les scripts du TP3 forment un mini-runbook.
\item \textbf{Pourquoi éviter d'ouvrir SSH à tout Internet ?}\\ Exposer le port 22 à \code{0.0.0.0/0} l'expose aux \textbf{scans} et attaques par \textbf{force brute}. On restreint à l'IP admin (\code{/32}).
\item \textbf{Quelle commande liste les ressources d'un groupe ?}\\ \code{az resource list --resource-group <rg>}.
\item \textbf{Quel service suit les métriques et crée des alertes ?}\\ \textbf{Azure Monitor} (métriques, logs, alertes, action groups).
\item \textbf{Différence entre métrique et log ?}\\ Une \textbf{métrique} est une valeur \textbf{numérique} dans le temps (CPU \%, débit) ; un \textbf{log} est un \textbf{événement détaillé} (texte structuré). La métrique mesure, le log raconte.
\item \textbf{Pourquoi versionner les scripts d'exploitation ?}\\ Pour \textbf{tracer} les modifications, revenir en arrière, collaborer (revue de code) et garantir la même version fiable. Un script d'exploitation est du code.
\item \textbf{Rôle de \code{DefaultAzureCredential} dans un script Python Azure ?}\\ Une \textbf{chaîne d'authentification} qui tente plusieurs sources (env, identité managée, \code{az login}) pour obtenir un token Azure AD \textbf{sans secret en dur}.
\item \textbf{Pourquoi tester un script sur un environnement non productif ?}\\ Un script peut \textbf{arrêter, modifier ou supprimer} des ressources. Le tester en dev évite d'impacter la production en cas de bug ou de mauvaise variable.
\item \textbf{Deux risques d'un script d'administration mal sécurisé ?}\\ \textbf{Action destructive sans confirmation} (arrêt/suppression accidentelle) ; \textbf{secrets en clair} dans le code (clé, mot de passe) exposés en cas de fuite.
\item \textbf{Quelle information doit figurer dans un rapport d'exploitation ?}\\ Inventaire, état des VM/services, contrôle des tags/gouvernance, sécurité, supervision, \textbf{analyse des coûts}, risques et \textbf{recommandations priorisées} — pour l'équipe technique et la DSI.
\item \textbf{Pourquoi un budget ne remplace-t-il pas une gouvernance des ressources ?}\\ Un budget \textbf{alerte sur un montant} mais ne dit pas quelles ressources, qui en est responsable, ni lesquelles arrêter. Sans tags ni organisation, on voit la dérive sans pouvoir agir ; il ne bloque pas la dépense.
\item \textbf{Deux actions simples pour réduire les coûts d'un environnement de dev ?}\\ \textbf{Désallouer les VM} hors usage (et détruire l'environnement entre séances) ; \textbf{supprimer les ressources orphelines / IP inutilisées} (ou réduire la taille des VM).
\item \textbf{Pourquoi journaliser les actions d'un script ?}\\ Pour \textbf{tracer} qui a fait quoi et quand (audit), diagnostiquer en cas d'incident et prouver les contrôles. \code{vm-power.sh} écrit dans \code{logs/vm-power.log}.
\item \textbf{Différence entre IaC et script d'exploitation ?}\\ L'\textbf{IaC} (Terraform) décrit l'\textbf{état cible} de façon déclarative (construit l'infra). Un \textbf{script d'exploitation} réalise des \textbf{actions ponctuelles} sur l'existant de façon impérative (administre l'infra). L'un construit, l'autre exploite.
\item \textbf{Citer trois contrôles d'un \code{healthcheck.sh}.}\\ Existence du \textbf{groupe de ressources} (bloquant) ; présence et \textbf{état des VM} ; \textbf{ressources sans tag} obligatoire (\code{Application}/\code{Owner}). \textit{(Également : alerte Azure Monitor, Storage + conteneur \code{operations}, règle NSG HTTP/HTTPS.)}
\end{enumerate}

\vspace{6mm}
\begin{center}
{\sffamily\bfseries\large\color{azuredark}\eParty~~Fin du dossier de rendu — TP3 Administration \& Automatisation Azure}\\[2mm]
{\color{black!55}ShopEasy — Louis \textsc{Scarfone} \& Maxence \textsc{Bourragué} — Bloc 4 · Cloud Computing}
\end{center}

\end{document}
